% ===================================================================
%  அட்டவணை எண் 13 — விடுதி. உணவகம். காப்பி விடுதி.
%
%  Traduction, faite en 2026, en tamoul d'aujourd'hui, sur l'ido de
%  texto/io. Regles de la colonne : voir 10-attavanai-01.tex. Recit a
%  la premiere personne, sans scenes. Les deux notes d'auteur sont en
%  gras-italique dans le fac-simile, contrairement a celles des
%  tableaux 1, 5 et 12 : la colonne suit. Ecrit avec outils/ordo.py
%  sous les yeux.
%
%  TRENTE-TROIS PAIRES POUR QUATRE-VINGT-SIX RENVOIS — le taux le
%  plus eleve de la colonne tamoule, et c'est le meme tableau qui
%  detient le record dans les colonnes marathe, telougoue et
%  coreenne. Quatre langues de quatre familles, quatre fois le meme
%  tableau en tete : la cause est dans le texte de 1926, pas dans les
%  langues d'arrivee. Le 13 ATTACHE — le conducteur DE l'omnibus, la
%  malle DU voyageur, la cuisine AU sous-sol, l'enseigne DE la halle,
%  la caissiere QUI inscrivait la facture SUR le registre.
%
%  LE MENU RESTE FRANCAIS, ET C'EST LA REGLE. Les tripes a la mode de
%  Caen, le gigot de mouton, le Saint-Emilion sont les CHOSES qu'on
%  mange dans cette salle-la ; les remplacer par des plats tamouls
%  aurait deplace le repas au lieu de traduire la langue. Le hindi,
%  le telougou et le coreen avaient tranche de meme.
%
%  ET LE SEUL JEU QUE LE TAMOUL N'EMPRUNTE PAS EST CELUI QUI EST
%  PARTI D'ICI. Le salon du bloc 12 aligne cinq jeux : le tamoul
%  ecrit டொமினோ, செக்கர்ஸ், திரிக்திராக், சீட்டாட்டம் — et
%  சதுரங்கம் pour les echecs (78). Ce mot-la n'est pas un emprunt :
%  c'est le nom sanskrit-tamoul du jeu des quatre corps d'armee, d'ou
%  viennent le shatranj persan, l'ajedrez espagnol et le chess
%  anglais. Sur une planche ou tout le reste du salon vient d'Europe,
%  un seul objet fait le chemin inverse, et la langue le montre sans
%  qu'on ait rien a expliquer.
%
%  LE JEU DE DAMES (79), LUI, N'A PAS DE NOM TAMOUL. Le pays a des
%  jeux de plateau anciens et nombreux — பல்லாங்குழி, ஆடு புலி
%  ஆட்டம், தாயம் —, et aucun n'est celui-la. La colonne telougoue
%  l'avait declare non resolu au meme renvoi ; celle-ci ecrit
%  l'emprunt செக்கர்ஸ், comme elle avait ecrit குரோக்கே et
%  பில்போக்கே au tableau 2. Coller un nom approchant ferait croire a
%  un jeu qu'on connait.
%
%  SIX ECARTS AU PREMIER PASSAGE, ET TOUS SUR DES PAIRES QUE
%  parigi.py LISTAIT. Ce n'est plus un angle mort, c'est le nombre :
%  trente-trois paires dans vingt-deux blocs, et la main finit par en
%  laisser passer. Le remede a chaque fois est le meme — nommer la
%  tete, rejeter le modificateur derriere un tiret. C'est une figure
%  juste, et c'est aussi une figure qu'on peut user : il faut la
%  compter. Ce tableau en emploie six pour vingt-deux blocs, ce qui
%  reste sous le taux du tableau 6 ; le jour ou elle deviendrait le
%  seul tour de la colonne, il faudrait varier, non parce que la
%  phrase serait fausse mais parce qu'elle serait de machine.
%
%  பால்கனி (bloc 10) POUR LA HUITIEME FOIS, et la regle n'a pas
%  bouge depuis le tableau 5.
% ===================================================================

%%K t13-noto-1 noto
\VUnotes{143.2mm}{%
\textit{\VUgras{(*) அறையின் விவரங்களுக்கு அட்டவணை எண் 6 ஐப் பார்க்க.}}%
}

%%K t13-apar-1 apar
\VUsaut{3.0mm}
\VUtitre{15.0pt}{60}{அட்டவணை எண் 13}
\VUsaut{1.6mm}
\VUpk{10.2pt}{40}{விடுதி. --- உணவகம்.}
\VUsaut{2.0mm}
\VUpk{10.2pt}{40}{காப்பி விடுதி.}
\VUsaut{3.4mm}

%%K t13-01-1 p
1. --- நேற்று மாலை ரோயானுக்கு வந்த பிறகு, « \textit{பெருங்கடல் விடுதி} »
என்னும் இடத்தில் தங்கினேன் ; அதை ஒரு நண்பர் எனக்குப் பரிந்துரைத்திருந்தார்.

%%K t13-01-2 p
அழகான \VUgras{அறை} \textsuperscript{(2)} ஒன்று எனக்குத் தரப்பட்டது —
இரண்டாம் \VUgras{தளத்தில்} \textsuperscript{(1)} ; அங்கே நான் மிக நன்றாகத்
தூங்கினேன் (*).

%%K t13-02-1 p
2. --- இன்று காலை, \VUgras{பணிப்பெண்} \textsuperscript{(3)} காலைச்
சிற்றுண்டியைக் கொண்டு வந்திருந்த என் அறையை விட்டு வெளியே வந்து,
\VUgras{புகைபிடிக்கும் அறைக்குள்} \textsuperscript{(5)} நுழைந்தேன் ; அது
\VUgras{வாசிப்பு அறையாகவும்} \textsuperscript{(4)} பயன்படுகிறது —
\VUgras{செய்தித்தாள்களைப்} \textsuperscript{(6)} படிப்பதற்காக,
\VUgras{சிகரெட்} \textsuperscript{(8)} புகைத்துக் கொண்டு. படித்த
பிறகு \VUgras{மின் தூக்கியில்} \textsuperscript{(9)} இறங்கி நகரத்தில்
உலாவப் போனேன்.

%%K t13-03-1 p
3. --- விடுதியின் \VUgras{எழுத்தரிடம்} \textsuperscript{(12)}
\VUgras{பொது மேசையில்} \textsuperscript{(11)} எந்த நேரத்தில் உணவு
தருகிறார்கள் என்று கேட்க மறந்திருந்ததால், முன்னதாகவே திரும்பி வந்து, அதைப் பயன்படுத்தி
விடுதியின் \VUgras{குளியலறையில்} \textsuperscript{(10)} குளிக்கப்
போனேன். பிறகு என் நண்பர்களுள் ஒருவருக்குத் தொலைபேச \VUgras{பணப் பெட்டி
அறைக்கு} \textsuperscript{(15)} மீண்டும் போனேன். ஆனால் \VUgras{தொலைபேசி}
\textsuperscript{(13)} பயன்பாட்டில் இருந்தது. என் திண்டாட்டத்தைப் பார்த்த
\VUgras{காசாளர்} \textsuperscript{(14)} — அவள் ஒரு \VUgras{கணக்குச்
சீட்டை} \textsuperscript{(17)} \VUgras{பயணிகள் பதிவேட்டில்}
\textsuperscript{(16)} பதிவு செய்து கொண்டிருந்தாள் —
\VUgras{வாயிற்காப்பாளரிடம்} \textsuperscript{(18)} கேட்டுக் கொண்டாள் ;
\VUgras{சிறு பணியாளை} \textsuperscript{(19)} \VUgras{மிதிவண்டியில்}
\textsuperscript{(20)} அனுப்பி என் செய்தியைச் சொல்லச் செய்யும்படி.

%%K t13-04-1 p
4. --- அவளுக்கு நன்றி சொல்லி, \VUgras{சுமை தூக்குபவருக்குச்}
\textsuperscript{(24)} சன்மானம் தர வெளியே வந்தேன் ; அவர் நிலையத்திலிருந்து
தன் \VUgras{கைவண்டியில்} \textsuperscript{(25)} என் \VUgras{தொப்பிப்
பேழையையும்} \textsuperscript{(26)}, \VUgras{என் கைப் பெட்டியையும்}
\textsuperscript{(27)}, \VUgras{என் பயணப் பையையும்}
\textsuperscript{(28)} கொண்டு வந்திருந்தார். இப்போதுதான் வந்த
\VUgras{கடிதம் விநியோகிப்பவர்} \textsuperscript{(21)} எனக்கு ஒரு
\VUgras{கடிதம்} \textsuperscript{(22)} தந்தார்.

%%K t13-05-1 p
5. --- \VUgras{அஞ்சல் பெட்டிக்கு} \textsuperscript{(23)} அருகில், கதவின்
வாசற்படியில் இன்னும் சிறிது நேரம் நின்று தெருவின் நடமாட்டத்தைப்
பார்த்தேன். \VUgras{ஓட்டுநர்} \textsuperscript{(30)} — \VUgras{பொது வண்டியின்}
\textsuperscript{(29)} ஓட்டுநர் — தன் வண்டியில் ஒரு \VUgras{பெட்டகத்தை}
\textsuperscript{(31)} ஏற்றிக் கொண்டிருந்தார் ; அது ஒரு
\VUgras{பயணியினுடையது} \textsuperscript{(32)} ; அந்தப் பயணியை
\VUgras{விடுதி உரிமையாளர்} \textsuperscript{(33)} வழியனுப்பிக்
கொண்டிருந்தார். இளம் \VUgras{தந்தி ஊழியன்} \textsuperscript{(35)}
அவசரம் இல்லாமல் விடுதியின் ஒரு வாடிக்கையாளருக்கு \VUgras{தந்தி}
\textsuperscript{(34)} ஒன்றைக் கொண்டு வந்தான் ; ஒரு \VUgras{சுற்றுலாப்
பயணி} \textsuperscript{(36)} தோளில் தன் \VUgras{ஒளிப்படக் கருவியுடன்}
\textsuperscript{(38)} தன் \VUgras{வழிகாட்டி நூலைப்}
\textsuperscript{(37)} பார்த்துக் கொண்டிருந்தார்.

%%K t13-tit-1 sub
\VUsaut{4.0mm}
\VUpk{10.2pt}{40}{மதிய உணவு நேரம்.}
\VUsaut{3.0mm}

%%K t13-06-1 p
6. --- கிராதிக்கு முன்னால் ஒரு கவலையற்ற \VUgras{தூது ஆள்}
\textsuperscript{(39)} தன் \VUgras{சுமை கொக்கிக்கும்} \textsuperscript{(40)}
தன் \VUgras{மெருகுப் பெட்டிக்கும்} \textsuperscript{(41)} இடையே
தூங்கிக் கொண்டிருந்தான் ; அதே வேளையில் இளம் \VUgras{துணி
துவைப்பவள்} \textsuperscript{(42)} — அவள் துணிகள் நிறைந்த ஒரு
\VUgras{கூடையைச்} \textsuperscript{(43)} சுமந்து வந்தாள் —
கதவருகில் நின்ற \VUgras{மேசைப் பணி அதிகாரியின்} \textsuperscript{(44)}
கம்பீரமான தோற்றத்தைப் பார்த்துச் சிரித்தாள்.

%%K t13-07-1 p
7. --- \VUgras{சமையல்காரன்} \textsuperscript{(45)} தன் \VUgras{சமையலறையை}
\textsuperscript{(47)} விட்டு வெளியே வருவதைப் பார்த்தேன் ; அந்தச்
சமையலறை \VUgras{நிலவறைத் தளத்தில்} \textsuperscript{(48)} இருக்கிறது ;
ஒரு \VUgras{பாத்திரம் கழுவுபவனை} \textsuperscript{(46)} செய்தி சொல்ல
அனுப்புவதற்காக அவன் வெளியே வந்தான். அதனால் உணவின் நேரம் நெருங்கி
விட்டது என்பது எனக்கு நினைவுக்கு வந்தது ; உணவகத்திற்குள் நுழைந்தேன் —
முற்றத்தைக் கடந்து ; அங்கே ஒரு \VUgras{மகிழுந்து ஓட்டுநர்}
\textsuperscript{(50)} தன் \VUgras{மகிழுந்தை} \textsuperscript{(49)}
நிறுத்தினார் ; அதே வேளையில் ஒரு \VUgras{இயந்திரப் பணியாளர்}
\textsuperscript{(52)} ஒரு \VUgras{மிதிவண்டிக்காரரின்}
\textsuperscript{(53)} அறிவுரைப்படி இப்போதுதான் விபத்துக்கு உள்ளான
\VUgras{எண்ணெய் முச்சக்கர வண்டியைப்} \textsuperscript{(51)} பழுது
பார்த்தார்.

%%K t13-08-1 p
8. --- ஒரு இளம் \VUgras{சிறு பணியாளிடம்} \textsuperscript{(54)} —
அவன் \VUgras{பீர்க் குவளைகளைச்} \textsuperscript{(55)} சுமந்து வந்தான் —
பொது மேசை
வராந்தாவின் கீழ் இருக்கிறதா என்று கேட்டேன்.

%%K t13-noto-2 noto
\VUnotes{143.2mm}{%
\textit{\VUgras{(*) அகராதியில் சேர்க்கப்பட்ட உணவுப் பட்டியலின்
துணையால், கீழே உள்ள உணவுப் பட்டியலை ஆசிரியர் எளிதாக மாற்றி
அமைக்கலாம்.}}%
}

%%K t13-08-1 p suite
அவன் ஆம் என்று பதில் சொன்னதால், மேசைகளுள் ஒன்றில் இடம் பிடித்து
மிகச் சிறந்த உணவைக் கொண்டு வரச் செய்தேன்.

%%K t13-tit-2 sub
\VUsaut{4.0mm}
\VUpk{10.2pt}{40}{மிகச் சிறந்த உணவு.}
\VUsaut{3.0mm}

%%K t13-09-1 p
9. --- பணியாளர் எனக்கு உணவின் பட்டியலையும் (*) மதுப் பட்டியலையும்
கொண்டு வந்தார். நான் தேர்ந்தெடுத்து ஆணையிட்டேன் : \VUgras{துணை
உணவாக} \VUgras{வெண்ணெயுடன்} \VUgras{இறால்கள்} ; தொடக்க உணவாக
\VUgras{கான் முறையில் குடல் உணவு}. \VUgras{மீன்} சாப்பிடவில்லை ;
ஆனால் ஆட்டுக் குட்டியின் \VUgras{தொடையில்} ஒரு பங்கு கேட்டேன் ;
\VUgras{காய்கறி உணவாக} ஏடுடன் \VUgras{கீரை}. பிறகு
\VUgras{இடை உணவுகளுள்} எதுவும் எனக்குப் பிடிக்காததால், பல
இனிப்புகளைக் கொண்டு வரச் செய்தேன் : \VUgras{சீஸ்},
\VUgras{பழங்கள்}, \VUgras{பிஸ்கட்டுகள்}. மேலும் மிகச் சிறந்த
சேன்-எமிலியோன் \VUgras{மதுவையும்} சேர்த்தேன் ; என் உணவைப் பற்றி
மிகவும் மனநிறைவு அடைந்து, \VUgras{பணியாளருக்கு} \textsuperscript{(57)}
நல்ல \VUgras{சன்மானம்} \textsuperscript{(59)} கொடுத்த பிறகு மேசையை
விட்டு எழுந்தேன்.

%%K t13-10-1 p
10. --- நான் புறப்படத் தயாராக இருப்பதைப் பார்த்த \VUgras{மேலாளர்}
\textsuperscript{(60)}, \VUgras{பெருங்கடலின்} \textsuperscript{(62)}
சிறந்த காட்சியை வியந்து பார்க்கப் பால்கனிக்குப் போகும்படி எனக்கு
முன்மொழிந்தார். தொலைவில் மீன்பிடி \VUgras{சிறு படகுகளும்}
\textsuperscript{(63)}, \VUgras{பாய்மரக் கப்பல்களும்}
\textsuperscript{(64)}, மாபெரும் \VUgras{அஞ்சல் கப்பலும்}
\textsuperscript{(65)} தெரியத் தொடங்கின ; அந்தக் கப்பலின் கருமையான
உருவம் வெண்மையான \VUgras{மேகங்களின்} \textsuperscript{(67)} மேல்
துருத்தித் தெரிந்தது — \VUgras{அடிவானத்தின்} \textsuperscript{(66)}
மேகங்கள். இலேசான தென்றல் \VUgras{கொடியை} \textsuperscript{(70)}
அசைத்தது ; அது \VUgras{பெயர்ப் பலகைக்கு} \textsuperscript{(69)} மேலே
நாட்டப்பட்டிருந்தது — \VUgras{கூடத்தின்} \textsuperscript{(68)}
பெயர்ப் பலகை.

%%K t13-tit-3 sub
\VUsaut{3.0mm}
\VUpk{10.2pt}{40}{பொழுதுபோக்குகள்.}
\VUsaut{3.0mm}

%%K t13-11-1 p
11. --- அந்தக் காட்சி அவ்வளவு சுவாரசியமாக இருந்ததால், நாளை
\VUgras{இரட்டைக் குழல் தொலைநோக்கியையோ} \textsuperscript{(71)}
\VUgras{தொலைநோக்கியையோ} \textsuperscript{(72)} கொண்டு காப்பி விடுதியின்
மொட்டை மாடத்தின் மேல் ஏற முடிவு செய்தேன்.

%%K t13-12-1 p
12. --- பால்கனியை விட்டு விடுதியின் காப்பி விடுதிக்குப் போனேன் ; ஒரு
\VUgras{பானத்தை} \textsuperscript{(73)} அருந்திக் கொண்டு \VUgras{பில்லியர்ட்ஸ்
ஆட்டம்} \textsuperscript{(75)} ஆடுவதற்காக. காப்பிக் கூடத்தை நிற்காமல்
கடந்தேன் ; அங்கே பல வாடிக்கையாளர்கள் \VUgras{சதுரங்கம்}
\textsuperscript{(78)}, \VUgras{செக்கர்ஸ்} \textsuperscript{(79)},
\VUgras{டொமினோ கற்கள்} \textsuperscript{(80)}, \VUgras{திரிக்திராக்}
\textsuperscript{(81)}, \VUgras{சீட்டாட்டம்} \textsuperscript{(82)}
ஆடிக் கொண்டிருந்தார்கள் ; நேராக \VUgras{பில்லியர்ட்ஸ் அறைக்கு}
\textsuperscript{(74)} ஏறினேன்.

%%K t13-13-1 p
13. --- ஆட்டம் முடிந்த போது தரைத் தளத்திற்கு இறங்கி, கடிதம் எழுத
பணியாளரிடம் \VUgras{உறை} \textsuperscript{(84)}, \VUgras{காகிதம்}
\textsuperscript{(83)}, \VUgras{அஞ்சல் தலை} \textsuperscript{(85)},
\VUgras{பேனா} \textsuperscript{(87)}, \VUgras{மை} \textsuperscript{(86)}
கேட்டேன்.

%%K t13-13-2 p
பிறகு ஒரு \VUgras{வண்டிக்காரரை} \textsuperscript{(92)} அழைத்தேன் ;
அவருடைய \VUgras{வண்டி} \textsuperscript{(88)} காப்பி விடுதிக்கு முன்னால்
நின்றிருந்தது. மணிக்கணக்கிலோ ஒரு ஓட்டத்திற்கோ விலை என்ன என்று
அவரிடம் கேட்டேன் ; விலை ஒத்துப் போன போது அவர் எனக்கு \VUgras{வண்டிக்
கதவைத்} \textsuperscript{(89)} திறந்து என்னை அமர வைத்தார். அவர்
மீண்டும் \VUgras{ஓட்டுநர் இருக்கையின்} \textsuperscript{(91)} மேல்
ஏறி, \VUgras{சிறு சாட்டையால்} \textsuperscript{(93)} குதிரைகளை
விரைவாகக் கிளப்பினார் ; அழகான உலாவுக்குப் புறப்பட்டோம்.

\VUsaut{6.0mm}
\VUfilet{22mm}
